\documentclass{beamer}
\usepackage[T1]{fontenc}
\usepackage[utf8]{inputenc}
\usepackage[spanish]{babel}
\usetheme{Madrid}
\title{Ayudantía 1 - Macroeconomía}
\author{Ignacio Garrido Urra}
\date{\textbf{Semestre 2025-2}}

\begin{document}
\frame{\titlepage}

\begin{frame}{Ejercicio 1}
Señale las diferencias entre \textbf{PIB} y el \textbf{Ingreso Disponible} ($Y_d$).
\end{frame}

\begin{frame}{Ejercicio 2}
Diga por qué no se contabiliza dentro del PIB a la compra y venta de \textit{insumos}.
\end{frame}

\begin{frame}{Ejercicio 3}
Considere la producción de un cierto país en donde se muestran sus respectivos bienes y precios en la tabla adjunta. 
Los bienes de los sectores \textbf{A y B} son vendidos a los hogares, los bienes del sector \textbf{C} son adquiridos por el gobierno, 
los bienes del sector \textbf{D} son exportados, mientras que los bienes del sector \textbf{E y F} son utilizados para fabricar los bienes A, B, C y D.

\medskip
\begin{center}
\begin{tabular}{lcccccc}
\hline
 & $P_A$ & $Q_A$ & $P_B$ & $Q_B$ & $P_C$ & $Q_C$ \\
\hline
2020 & \$500 & 100 & \$10 & 2000 & \$100 & 80 \\
2021 & \$540 & 120 & \$9 & 1800 & \$120 & 90 \\
\hline
\end{tabular}

\medskip
\begin{tabular}{lcccccc}
\hline
 & $P_D$ & $Q_D$ & $P_E$ & $Q_E$ & $P_F$ & $Q_F$ \\
\hline
2020 & \$10 & 600 & \$100 & 700 & \$20 & 1000 \\
2021 & \$15 & 500 & \$80 & 900 & \$30 & 1100 \\
\hline
\end{tabular}
\end{center}
\end{frame}


\begin{frame}{Ejercicio 3}
\medskip
\begin{enumerate}
\item[(a)] Determine el \textit{PIB nominal} y \textit{real} para cada año.
\item[(b)] Calcule la \textit{tasa de crecimiento} del PIB real usando los precios del año 2020 como base.
\item[(c)] Calcule el \textit{deflactor del PIB}.
\item[(d)] Calcule el \textit{IPC} para cada período con bienes A y B.
\item[(e)] Calcule la \textit{inflación} por deflactor e IPC.
\end{enumerate}
\end{frame}

\begin{frame}{Ejercicio 4 (Certamen 2015)}
Explique si en una economía cerrada (sin exportaciones ni importaciones), en la cual las demandas por bienes y servicios finales son inelásticas al precio, 
el \textbf{costo de vida} medido por el \textit{IPC} o por el \textit{deflactor del PIB} deberían ser iguales.
\end{frame}

\begin{frame}{Ejercicio 5 (Certamen 2013)}
En el país de los jóvenes se consumen solo cuatro artículos (Ron, cocacola, educación y PlayStation). 
En el año 2010 una familia representativa consumía 3 unidades de ron, 3 unidades de cocacola, 2 unidades de educación y 1 unidad de PlayStation. 
Los precios (en \$) durante los últimos años son:

\begin{center}
\begin{tabular}{lcccc}
\hline
 & 2010 & 2011 & 2012 & 2013 \\
\hline
Ron & 400 & 350 & 450 & 500 \\
Cocacola & 1000 & 1100 & 1100 & 1150 \\
Educación & 600 & 800 & 1000 & 1020 \\
PlayStation & 800 & 1200 & 1200 & 1300 \\
\hline
\end{tabular}
\end{center}

\medskip
Calcule el \textbf{IPC} (año base 2010) y la \textbf{inflación} en 2011, 2012 y 2013.
\end{frame}

\begin{frame}{Ejercicio 6 (Certamen 2013)}
Datos de cuentas nacionales (en miles):

\begin{center}
\begin{itemize}
\item PIB = 60.000
\item Inversión Bruta = 12.000
\item Inversión Neta = 9.000
\item Consumo = 33.000
\item Gasto del Gobierno = 14.500
\item Superávit Fiscal = 500
\end{itemize}
\end{center}
\medskip
Calcule el \textbf{PIN}, las \textbf{Exportaciones Netas}, los \textbf{impuestos menos transferencias}, el \textbf{Ingreso Disponible} y el \textbf{Ahorro}.
\end{frame}


\begin{frame}{Reflexionemos (Noticia Emol)}
\centering
\includegraphics[width=\textwidth,height=0.72\textheight,keepaspectratio]{1Ayudantia/noticia1.PNG}
\medskip

PIB de Chile creció un 2,3\% anual 
\href{https://www.emol.com/noticias/Economia/2025/05/19/1166723/crecimiento-primer-trimestre.html}{(Emol, Mayo 2025)}.


\end{frame}



\begin{frame}{Resumen Noticia: Actualidad Económica (Primer Trimestre 2025)}
PIB de Chile creció un \textbf{2,3 \%} anual, superando la expectativa de 2 \%.  Este aumento fue explicado mayormente por el impulso de:
\begin{itemize}
  \item Exportaciones (+10,7 \%) e importaciones (+9,0 \%) con impacto neto positivo en el PIB.
  \item Mayor consumo de los hogares (+1,8 \%), especialmente bienes durables (tecnología).
  \item Gasto del gobierno (+3,1 \%), liderado por inversión en salud.
  \item Actividades que destacaron: comercio, manufactura, servicios personales y agro-silvícola.  
\end{itemize}
\medskip
\textit{\footnotesize Fuente: Banco Central vía Emol — 19 de mayo de 2025.}
\end{frame}

\begin{frame}{Reflexión en base a la noticia}
\textbf{PIB creció 2,3 \% en 1T 2025, superando expectativas. Impulso: exportaciones + consumo + gasto público.}

\medskip
\textbf{Preguntas para discutir:}
\begin{enumerate}
  \item ¿Qué componente (exportaciones, consumo, gasto público) explicaría mejor el crecimiento real? ¿Por qué?
  \item Según nuestros conceptos, ¿pierde valor o gana relevancia el PIB nominal frente al real en este escenario?
  \item ¿El costo de vida de los hogares cambió de forma similar a lo que indica el deflactor del PIB?
  \item Si hay aumento en consumo de durables, ¿qué efectos tendría en la inversión en el largo plazo?
  \item ¿Qué vulnerabilidades identificas ante un crecimiento tan dependiente de las exportaciones?
\end{enumerate}
\end{frame}


\end{document}
